\documentclass{book}

\usepackage{amsmath}
\usepackage{amsfonts}
\usepackage{amssymb}
\usepackage{amscd}
\usepackage{amsthm}
\usepackage{mathrsfs}
\usepackage{graphics}
\usepackage{graphicx}
\usepackage{tikz}
\usepackage{tikz-cd}
\usepackage{bm}

\newtheorem{theorem}{Theorem}
\newtheorem{definition}{Definition}
\newcommand{\bracenom}{\genfrac{\lbrace}{\rbrace}{0pt}{}}

\graphicspath{ {./} }

\title{Category Theory Leinster - Solutions Manual}
\author{Mingruifu Lin}
\date{September 2026}

\begin{document}

\maketitle

\tableofcontents

\chapter{Categories, functors and natural Transformations}

\section{Categories}

\subsection*{Exercise 1.1.12}
I'm not someone with much ideas.

\subsection*{Exercise 1.1.13}
Suppose both $g_1$ and $g_2$ are inverses of $f$. First, we have
$$g_1f = 1_A$$
Multiply both sides by $g_2$
$$g_1fg_2 = 1_A g_2$$
Since $fg_2 = 1_B$, we have
$$g_1 1_B = 1_A g_2$$
Note that the difference in identity makes sense ($g_1$ requires type $B$, but $g_2$ produces type $A$). Anyways, applying identity, we end up with
$$g_1 = g_2$$

\subsection*{Exercise 1.1.14}
Composition is defined
$$\mathscr{A}((A', B'), (A'', B'')) \times \mathscr{A}((A, B), (A', B')) \rightarrow \mathscr{A}((A, B), (A'', B''))$$
$$(g_A \times g_B) \times (f_A \times f_B) \mapsto (g_A \circ f_A) \times (g_B \circ f_B)$$
And identity is defined
$$1_{(A, B)} = 1_A \times 1_B$$

\subsection*{Exercise 1.1.15}
I don't understand the question, nor do I know what's a homotopy class.

\end{document}